\chapter{前言}
概率论从给定模型推导随机结果，数理统计从有限样本反推模型中的未知参数。
复习的关键是识别“随机什么、已知什么、在哪个条件下计算”。
本册把事件、分布、数字特征、极限与统计推断串成一条主线，计算旁保留条件，结论旁给出证明或可追踪的证明思路。

第1—8章覆盖考研常见核心知识；“考研核心”是学习层级标记，不声称替代某一年度、某一科目的官方考试大纲。
第9章安排方差分析与一元回归初步；第10章独立介绍测度论语言，可在完成主线后阅读。
竞赛风格题强调计数、对称、指示变量、条件分解、构造反例和多方法互证。

本册以标准概率模型、统计推断原理和可复现的数学推导组织知识。正文引用用于定位知识主题；未标为原题的训练均为本册编写的风格题。补充理论保留适用条件，重要内容配有理解札记、易错辨析和数形结合图示。

全书150个例题分别给出题目、思路、解答、点评和易错提醒；80道章末习题与答案双向链接。
多方法题的目的在于训练选择：补事件减少分类，条件化降低维数，变换揭示支撑，正交投影解释自由度。
读者可先独立解题，再对照方法点评，最后完成不同参数或不同假设下的变式。

所有图形由TikZ或pgfplots绘制，数学对象与正文同步。数字近似保留明确的近似号，精确有限样本结论与渐近结论分开标识。
